% Chapter 2

\chapter{Revisão da Literatura}	%The main chapter title
\label{Chapter2}

Após a Introdução \ref{Chapter1} onde foi definida a metodologia \ref{sec:metodologia}, agora serão apresentados  e analisados os principais artigos selecionados e os respetivas fontes e autores mais relevantes sobre o tema em estudo.
Neste capítulo após a análise bibliográfica será feita uma apreciação crítica dos principais artigos selecionados e as suas contribuições para o tema em estudo.
%1. Pesquisa Bibliográfica (c/ base numa análise bibliométrica c/ mínimo de 30 referências bibliográficas) • Definir a estratégia de pesquisa; • Identificar as bases de dados relevantes; • Definir os critérios de inclusão/exclusão.

Dos 536 artigos fornecidos pelo \gls{wos} através da filtragem feita anteriormente,foram destacados os seguintes:
\textbf{Artigo I:}  "A Review of Technical Standards for Smart Cities"\cite{article}.
\textbf{Artigo II:} "AI-Driven Optimization of Urban Logistics in Smart Cities: Integrating Autonomous Vehicles and IoT for Efficient Delivery Systems"\cite{su162411265}.
\textbf{Artigo III:} "Collaborative urban transportation: Recent advances in theory and practice"\cite{CLEOPHAS2019801}.
\textbf{Artigo IV:} "Designing new models for energy efficiency in urban freight transport for smart cities and its application to the Spanish case"\cite{NAVARRO2016314}.
\textbf{Artigo V:} "Fulfilment of last-mile urban logistics for sustainable and inclusive smart cities: a case study conducted in Portugal"\cite{Correia02062024}.
\textbf{Artigo VI:} "Interaction with City Logistics Stakeholders as a Factor of the Development of Polish Cities on the Way to Becoming Smart Cities"\cite{en15114103}.
\textbf{Artigo VII:} "Shortening the Last Mile in Urban Areas: Optimizing a Smart Logistics Concept for E-Grocery Operations"\cite{GonzalezRomero06102024}.

%3. Análise bibliométrica + síntese crítica

%Como funciona: Primeiro apresenta a análise bibliométrica (mapas, gráficos, autores mais citados),
%depois faz uma síntese crítica dos principais artigos (não necessariamente um por seção, mas destacando os 7 mais relevantes dentro dos temas).

%%%%%%%%%%%%%%%%%%%%%%%%%%%%%%%%%%%%%%%%%%%%%%%%%%%%%%%%%%%%%%%%%%%%%%%%%%%%%%%
\section{Análise Bibliométrica}
\label{sec:analise_bibliometrica}

%2. Revisão da Literatura (c/ base numa análise bibliométrica c/ mínimo de 30 referências bibliográficas) • Identificar os principais contributos e seus autores; • Identificar áreas de influência e interesse. Maior ênfase em 7 referências bibliográficas;Peso maior!
A análise Bibliométrica aqui presente é um fator crucial para identificar os principais artigos, fontes e autores filtrados de acordo com os artigos escolhidos, esta foi criada em um ficheiro \gls{excel} submetido para avaliação e que está disponível no repositório \gls{github} associado a este relatório \cite{JunqueiraDevEduardoRepositório}.
A revisão da literatura foca-se apenas nestes 7 artigos selecionados sendo que os restantes 529 artigos serviram apenas para contextualização de mapas no \gls{vosviewer}, como já foi referido anteriormente.

Assim, esta análise têm como objetivo identificar os principais indicadores de desempenho \gls{kpi}.

Para tal foram calculados os seguintes indicadores para os \textbf{Artigos} selecionados:
\begin{itemize}
    \item Ano de publicação - com maior ênfase nos últimos 10 anos (2016 a 2026);
    \item nº de citações nos artigos tendo como base a plataforma do \gls{wos} variando de (17 a 224 citações);
    \item Fontes de publicação mais relevantes com a \gls{query};
    \item nº de citação do artigo por ano (4 a 32 citações/ano);
    \item nº de referências no artigo (6 a 155 referências);
    \item autores mais relevantes (quantidade de vezes que o autor é citado em artigos da sua posse);
    \item nº de autores em cada artigo (1 a 10) .
    \item \textit{Keywords} do artigo mais relevantes com a \gls{query}.
    \item Interpretação do \textit{Abstract} que seja relevante.
    \item Países como Portugal, Espanha e Polónia que são filtros de escolha adicionais por Eduardo Junqueira. 
\end{itemize}

Agora os seguintes indicadores para as \textbf{Fontes} destes artigos anteriores:
\begin{itemize}
    \item Fonte de publicação: (Clean Technologies; Sustainability
European Journal of Operational Research (EJOR);
Transportation Research Procedia;
International Journal of Logistics Research and Applications;
Energies;
Smart Cities.
);
    \item Nível dos últimos 5 anos de impacto do Jornal (3.60 a 6.20);
    \item Categoria do Quartil do Jornal (Q1 a Q3);
    \item Categoria rank do Jornal (14 a 112);
 \end{itemize}

Por fim  os seguintes indicadores para os \textbf{Autores} destes artigos:
\begin{itemize}
    \item Autor mais relevante consoante todos os patâmetros abaixo;
    \item nº de publicações do autor na plataforma do \gls{wos} variando de (3 a 303 publicações);
    \item nº de citação do autor desde sempre na plataforma do \gls{wos} (19 a 40957 citações);
    \item nº de referências no artigo (6 a 155 referências);
    \item classificação flobal do autor \textit{H-index} na plataforma do \gls{wos} (1 a 41);
\end{itemize}

%%%%%%%%%%%%%%%%%%%%%%%%%%%%%%%%%%%%%%%%%%%%%%%%%%%%%%%%%%%%%%%%%%%%%%%%%%%%%%%%
\section{Apreciação Crítica dos Principais Artigos}
\label{sec:apreciacao_critica_principais_artigos}


